% Part: computability
% Chapter: recursive-functions
% Section: other-recursions

\documentclass[../../../include/open-logic-section]{subfiles}

\begin{document}

\olfileid{cmp}{rec}{ore}
\olsection{عوديات أخرى}

وبالاقتران والجمع في متتاليات نرد إلى العودية البدائية صورًا أخرى،
أبعد عن صورتها الأولى وأوسع نفعًا. فمن ذلك تعريف دالتين معًا،
على الوجه الآتي:
\begin{align*}
h_0(\vec x, 0) & = f_0(\vec x) \\
h_1(\vec x, 0) & = f_1(\vec x) \\
h_0(\vec x, y+1) & = g_0(\vec x, y, h_0(\vec x, y), h_1(\vec x, y)) \\
h_1(\vec x, y+1) & = g_1(\vec x, y, h_0(\vec x, y), h_1(\vec x, y))
\end{align*}
وهذا ما يسمى \emph{العودية المتزامنة}. ومن وجوه التعريف النافعة أيضًا
أن نعيّن $h(\vec x,y+1)$ بدلالة \emph{جميع} القيم
$h(\vec x,0)$، …،~$h(\vec x,y)$، كما في التعريف الآتي:
\begin{align*}
h(\vec x, 0) & = f(\vec x) \\
h(\vec x, y+1) & = g(\vec x, y, \tuple{h(\vec x, 0), \dots, h(\vec x, y)}).
\end{align*}
ويجمع المخطط الآتي هذا المعنى في عبارة أقصر:
\[
h(\vec x, y) = g(\vec x, y, \tuple{h(\vec x, 0), \dots, h(\vec x, y-1)})
\]
على أن تكون الحجة الأخيرة لـ$g$ المتتالية الخالية عند
$y=0$. والمعنى في الصورتين واحد: للدالة $h$ عند حساب
«خطوة اللاحق» أن تستعمل سائر القيم التي حُسبت قبلها، مجموعة في متتالية.
وهذا هو \emph{عودية مسار القيم}. ومن تطبيقاته
جواز تعريف من القبيل الآتي:
\begin{align*}
h(\vec x, y) & = \begin{cases}
  g(\vec x, y, h(\vec x, k(\vec x, y))) & \text{إذا كان $k(\vec x, y) < y$} \\
  f(\vec x) & \text{في غير ذلك}
\end{cases}
\end{align*}
فتُحسب قيمة $h$ عند $y$ من قيمة $h$ عند
\emph{أي} حجة سابقة، تتولى~$k$ تعيينها.

\begin{prob}
  ضع لدالة الباقي $r(x,y)$ تعريفًا بعودية مسار القيم. (إذا كان $x$ و$y$
  عددين طبيعيين وكان $y>0$، فإن $r(x,y)$ هو العدد الأصغر من~$y$ الذي
  يحقق $x=z\times y+r(x,y)$ لبعض~$z$. وللتعيين، لنقل إنه إذا كان
  $y=0$ فإن $r(x,0)=0$.)
\end{prob}

فتأمل كيف تُستخرج هذه الدوال بالعودية البدائية
المعتادة. وبقيت صورة أوسع تصرفًا، نجيز فيها
تبديل \emph{المعلمات}، أي القيم الجانبية، أثناء العودية:
\begin{align*}
h(\vec x, 0) & = f(\vec x) \\
h(\vec x, y+1) & = g(\vec x, y, h(k(\vec x), y))
\end{align*}
وهذه أيضًا تحاكيها العودية البدائية المعتادة. وردّها إليها يحتاج إلى دقة؛
ومن سبيل الاهتداء إليه أن تبسط الحساب بيدك خطوة خطوة.

\end{document}
